\begin{lemma}[Uniform discriminant estimate]\label{lem:uniform-quadratic-discriminant}
Let $n\geq1$, and write $9n^2+4=t^2d$ with $t\geq1$ and $d>1$ squarefree.
For the conductor $f$ of $\mathbb Q(\sqrt d)$, put
$\lambda=(\varphi(f)-|\mu(f)|)/\sqrt{f/d}$ as in
\cite[Lemma~\ref{B-lem:odd-quadratic-length}]{Finite}. Then
\[
 \frac{\lambda}{\sqrt d}\geq\frac8{\sqrt{10}}.
\]
\end{lemma}
\begin{proof}
The integer $9n^2+4$ lies strictly between $(3n)^2$ and $(3n+1)^2$,
so $d>1$. Every odd prime $p\mid d$ satisfies $p\ne3$ and
$(3n)^2\equiv-4\pmod p$. Therefore $p\equiv1\pmod4$; in particular,
$p\geq5$. Reducing $t^2d=9n^2+4$ modulo three gives $d\equiv1\pmod3$.

First suppose $d$ is odd. Then $d\equiv1\pmod4$, $f=d$, and
$\lambda=\varphi(d)-1$. Since $d\ne1,5$, it has a prime factor $p\geq13$.
Write $d=ph$. Every prime factor $q$ of $h$ satisfies $q\geq5$, whence
$(q-1)/\sqrt q\geq1$ and $\varphi(h)\geq\sqrt h$. Consequently
\[
 \frac{\varphi(d)-1}{\sqrt d}
 =\frac{(p-1)\varphi(h)-1}{\sqrt{ph}}
 \geq\frac{(p-2)\varphi(h)}{\sqrt{ph}}
 \geq\frac{p-2}{\sqrt p}
 \geq\frac{11}{\sqrt{13}}>\frac8{\sqrt{10}}.
\]
Here $(p-2)/\sqrt p=\sqrt p-2/\sqrt p$ increases with $p>0$.

If $d$ is even, write $d=2h$. Then $h$ is odd, $h\geq5$ because
$d\equiv1\pmod3$ excludes $d=2$, and every prime factor of $h$ is
$1\pmod4$. Now $f=8h$, $\mu(f)=0$ and
$\lambda=2\varphi(h)$. Each factor $(p-1)/\sqrt p$ of
$\varphi(h)/\sqrt h$ is at least $4/\sqrt5>1$, and there is at least
one such factor. Hence
\[
 \frac{\lambda}{\sqrt d}
 =\sqrt2\frac{\varphi(h)}{\sqrt h}
 \geq\frac{4\sqrt2}{\sqrt5}=\frac8{\sqrt{10}}.
\]
\end{proof}

